%%%%%%%%%%%%%%%%%%%%%%%%%%%%%%%%%%%%%%%%%%%%%%%%%%%%%%%%%%%
\begin{frame}[fragile]\frametitle{}

	\begin{center}
	{\Large Challenges Now}
	\end{center}

\end{frame}


%%%%%%%%%%%%%%%%%%%%%%%%%%%%%%%%%%%%%%%%%%%%%%%%%%%%%%%%%%%
\begin{frame}[fragile]\frametitle{The Challenge Map: Who Is Struggling With What?}

\begin{center}
\begin{tabular}{|l|p{7cm}|}
\hline
\textbf{Role} & \textbf{Core Challenge} \\
\hline
Student (STEM) & Framework overload, tutorial hell, no deployed projects \\
\hline
Student (Non-STEM) & ``AI is not for me'' barrier, no clear entry point \\
\hline
Fresher (Engineering) & College curriculum 3--5 years behind, entry-level jobs need experience \\
\hline
Fresher (Non-Engineering) & Resume doesn't signal AI-readiness, competing with engineers \\
\hline
Lead / Architect & Their output is partly AI-automatable, system design patterns unclear \\
\hline
Manager & AI ROI pressure, risk of AI theater, workforce transition decisions \\
\hline
\end{tabular}
\end{center}

\vspace{0.2cm}
{\small Common thread: \textbf{speed of change outpaces any individual's ability to absorb it}}

\end{frame}


%%%%%%%%%%%%%%%%%%%%%%%%%%%%%%%%%%%%%%%%%%%%%%%%%%%%%%%%%%%
\begin{frame}[fragile]\frametitle{Students --- STEM Background}

\begin{columns}
    \begin{column}[T]{0.55\linewidth}
\textbf{What they face:}
\begin{itemize}
\item \textbf{Framework paralysis:} PyTorch or JAX? LangChain or LlamaIndex? A new one every month
\item \textbf{Tutorial hell:} Watching 50 hours of courses, building nothing deployable
\item \textbf{Wrong benchmark:} Over-indexed on Kaggle medals; industry now wants RAG and agents, not XGBoost tuning
\item \textbf{Imposter syndrome:} ``I need a PhD to work in real AI''
\item \textbf{Maths anxiety:} Trying to master backprop before writing a single API call
\end{itemize}
    \end{column}
    \begin{column}[T]{0.45\linewidth}
\textbf{The hidden danger:}

Skill half-life in AI is 3--4 months. A student who spends 2 years perfecting 2020-era skills graduates into a 2025 job market that has moved on.

\vspace{0.3cm}
\textbf{What employers actually say:}

``Show me one working app, not ten certificates.''
    \end{column}
\end{columns}

\end{frame}


%%%%%%%%%%%%%%%%%%%%%%%%%%%%%%%%%%%%%%%%%%%%%%%%%%%%%%%%%%%
\begin{frame}[fragile]\frametitle{Students: Non-STEM Background}

\begin{columns}
    \begin{column}[T]{0.55\linewidth}
\textbf{What they face:}
\begin{itemize}
\item \textbf{False belief:} ``AI is only for engineers and mathematicians''
\item \textbf{No visible entry point:} Every roadmap starts with Python, linear algebra, calculus, immediately alienating
\item \textbf{Domain expertise invisible:} A psychology or law student does not see their knowledge as an AI asset
\item \textbf{Under-representation:} AI communities and events skew heavily technical; non-STEM students feel out of place
\item \textbf{Credential confusion:} Should I do a PG Diploma? A bootcamp? Switch branches?
\end{itemize}
    \end{column}
    \begin{column}[T]{0.45\linewidth}
\textbf{The reality they don't see:}

\begin{itemize}
\item A psychology student understands human bias, critical for AI ethics and UX
\item A law student can evaluate AI contracts and compliance, a scarce and growing skill
\item A journalist can prompt-engineer better than most engineers
\end{itemize}

\vspace{0.2cm}
Domain expertise \textbf{is} the moat. The coding layer is thinner than ever.
    \end{column}
\end{columns}

\end{frame}


%%%%%%%%%%%%%%%%%%%%%%%%%%%%%%%%%%%%%%%%%%%%%%%%%%%%%%%%%%%
\begin{frame}[fragile]\frametitle{Freshers, Engineering / STEM Graduates}

\begin{columns}
    \begin{column}[T]{0.55\linewidth}
\textbf{What they face:}
\begin{itemize}
\item \textbf{Curriculum lag:} College syllabi are 3--5 years behind; still teaching scikit-learn basics while industry runs on LLM APIs and vector databases
\item \textbf{Experience paradox:} Entry-level job descriptions demand 2--3 years of AI experience
\item \textbf{No production thinking:} Projects run in Jupyter notebooks; no deployment, no monitoring, no cost awareness
\item \textbf{GenAI FOMO:} Jumping between ChatGPT wrappers without building depth in any one area
\item \textbf{Resume noise:} Hundreds of applicants with identical MOOC certificates
\end{itemize}
    \end{column}
    \begin{column}[T]{0.45\linewidth}
\textbf{What differentiates:}

\begin{itemize}
\item One deployed app with real users (even 10 users)
\item An open-source PR merged on HuggingFace or LangChain
\item A blog post explaining \textit{why} something worked, not just that it did
\end{itemize}

\vspace{0.2cm}
{\small Employers can tell in 5 minutes whether a candidate has \textbf{built} or only \textbf{watched}.}
    \end{column}
\end{columns}

\end{frame}


%%%%%%%%%%%%%%%%%%%%%%%%%%%%%%%%%%%%%%%%%%%%%%%%%%%%%%%%%%%
\begin{frame}[fragile]\frametitle{Freshers: Non-Engineering / Non-STEM Graduates}

\begin{columns}
    \begin{column}[T]{0.55\linewidth}
\textbf{What they face:}
\begin{itemize}
\item \textbf{Competing with engineers} who also want to ``move into AI'', and the market defaults to hiring them first
\item \textbf{Resume does not signal AI readiness:} A commerce or arts degree triggers an automatic filter in ATS (Applicant Tracking Systems)
\item \textbf{Credential dilemma:} MBA + AI specialization, or MS in AI? Two years and large fees either way
\item \textbf{Role confusion:} AI Product Manager? AI Consultant? AI Trainer? The new roles are poorly defined and hard to target
\item \textbf{Network gap:} Engineering peers have lab networks, project groups, hackathon teams
\end{itemize}
    \end{column}
    \begin{column}[T]{0.45\linewidth}
\textbf{The opening they miss:}

\begin{itemize}
\item AI content and prompt engineering roles explicitly \textit{prefer} language, communication, and domain skills over coding
\item AI ethics, policy, and compliance teams are growing and are nearly always led by non-engineers
\item ``AI Trainer'' roles at RLHF (Reinforcement Learning from Human Feedback) companies (Scale AI, Outlier) pay well and require zero coding
\end{itemize}
    \end{column}
\end{columns}

\end{frame}


%%%%%%%%%%%%%%%%%%%%%%%%%%%%%%%%%%%%%%%%%%%%%%%%%%%%%%%%%%%
\begin{frame}[fragile]\frametitle{Lead-Level Professionals (5--12 Years Experience)}

\begin{columns}
    \begin{column}[T]{0.55\linewidth}
\textbf{What they face:}
\begin{itemize}
\item \textbf{Output under threat:} AI tools now do 30--50\% of what a mid-level analyst, developer, or writer produces, faster
\item \textbf{Workflow inertia:} Decade of mastered habits; learning new tools feels slow and risky
\item \textbf{Direction confusion:} Prompt engineering? Fine-tuning? Agents? Where to invest limited learning time?
\item \textbf{Authority erosion:} Juniors using AI tools are producing outputs that rival senior work
\item \textbf{Visibility risk:} Being seen as the person who ``doesn't get AI'' by management
\end{itemize}
    \end{column}
    \begin{column}[T]{0.45\linewidth}
\textbf{The bigger risk:}

This group has the most to lose from \textit{inaction} and the most to gain from \textit{early adoption}.

\vspace{0.3cm}
A lead who learns to orchestrate AI becomes a \textbf{force multiplier}. A lead who ignores it becomes \textbf{redundant}.

\vspace{0.3cm}
{\small ``Your 10 years of domain judgment is irreplaceable. Add AI execution on top and you are unstoppable.''}
    \end{column}
\end{columns}

\end{frame}


%%%%%%%%%%%%%%%%%%%%%%%%%%%%%%%%%%%%%%%%%%%%%%%%%%%%%%%%%%%
\begin{frame}[fragile]\frametitle{Architects \& Senior Technical Leaders}

\begin{columns}
    \begin{column}[T]{0.55\linewidth}
\textbf{What they face:}
\begin{itemize}
\item \textbf{Non-determinism problem:} LLMs break traditional software design; the same input can yield different outputs, deterministic architecture patterns do not apply
\item \textbf{Evaluation gap:} How do you measure if an AI system is ``correct''? No unit test covers hallucination
\item \textbf{Cost \& latency:} Token costs, rate limits, and response times are design constraints unlike any before
\item \textbf{Fastest-moving infra in history:} The model you designed for last quarter may already be deprecated
\item \textbf{Responsibility without precedent:} Accountable for AI system failures with no established playbook
\end{itemize}
    \end{column}
    \begin{column}[T]{0.45\linewidth}
\textbf{New design concerns:}

\begin{itemize}
\item Retrieval-Augmented Generation (RAG) vs. fine-tuning vs. in-context learning, which tier for which problem?
\item Guardrails, output validation, fallback chains
\item Model versioning and prompt drift
\item Vendor lock-in risk (OpenAI today, what tomorrow?)
\end{itemize}

\vspace{0.2cm}
{\tiny (Ref: Chip Huyen, ``Designing ML Systems''; Martin Fowler, AI Architecture Patterns)}
    \end{column}
\end{columns}

\end{frame}


%%%%%%%%%%%%%%%%%%%%%%%%%%%%%%%%%%%%%%%%%%%%%%%%%%%%%%%%%%%
\begin{frame}[fragile]\frametitle{Managers \& Business Leaders}

\begin{columns}
    \begin{column}[T]{0.55\linewidth}
\textbf{What they face:}
\begin{itemize}
\item \textbf{ROI pressure from above:} Leadership demands AI results; timelines are unrealistic
\item \textbf{AI theater risk:} Buying ChatGPT Team subscriptions and calling it ``AI transformation'', no real change in output quality
\item \textbf{Prioritization paralysis:} 20 AI vendor demos a week, no framework to evaluate which matters
\item \textbf{Workforce transition:} Who to upskill, who to redeploy, who the team will lose to AI-automatable tasks
\item \textbf{Trust deficit:} Teams fear AI means layoffs; manager caught between leadership and team
\end{itemize}
    \end{column}
    \begin{column}[T]{0.45\linewidth}
\textbf{The core dilemma:}

\vspace{0.2cm}
Move too slow $\rightarrow$ competitors AI-transform first, you lose market share.

\vspace{0.2cm}
Move too fast $\rightarrow$ poor implementations, hallucinating AI in customer-facing roles, trust destroyed.

\vspace{0.3cm}
{\small ``The manager's job is no longer to manage people doing work. It is to manage people \textit{and AI} doing work together.''}
    \end{column}
\end{columns}

{\tiny (Ref: WEF Future of Jobs Report 2024; HBR ``Managing the AI Transition'' 2024)}
\end{frame}


%%%%%%%%%%%%%%%%%%%%%%%%%%%%%%%%%%%%%%%%%%%%%%%%%%%%%%%%%%%
\begin{frame}[fragile]\frametitle{}

\begin{center}
{\Large What Can Be Done?}
\end{center}

\end{frame}


%%%%%%%%%%%%%%%%%%%%%%%%%%%%%%%%%%%%%%%%%%%%%%%%%%%%%%%%%%%
\begin{frame}[fragile]\frametitle{For Students: Stop Watching, Start Building}

\begin{columns}
    \begin{column}[T]{0.5\linewidth}
\textbf{STEM students:}
\begin{itemize}
\item Skip framework debates, pick one (LangChain or LlamaIndex) and ship something
\item Build one deployed app using a free-tier API (OpenAI, Groq, Gemini)
\item Join AI hackathons: Smart India Hackathon, MLH, Devfolio
\item Contribute one pull request to a HuggingFace or LangChain repo
\item Read one paper/week from Arxiv Sanity or Papers With Code
\end{itemize}
    \end{column}
    \begin{column}[T]{0.5\linewidth}
\textbf{Non-STEM students:}
\begin{itemize}
\item Start with AI tools in your own domain: legal research, psychology surveys, journalism
\item Take ``AI for Everyone'' (Andrew Ng, Coursera), no coding required
\item Write one blog post explaining an AI concept in plain language
\item Attend one AI community event (Pune AI, TFUG, MLPC)
\item Build a portfolio of \textit{AI-augmented domain work}, not a coding portfolio
\end{itemize}
    \end{column}
\end{columns}

\vspace{0.2cm}
{\small Rule of thumb: \textbf{one deployed thing} is worth more than ten certificates.}

\end{frame}


%%%%%%%%%%%%%%%%%%%%%%%%%%%%%%%%%%%%%%%%%%%%%%%%%%%%%%%%%%%
\begin{frame}[fragile]\frametitle{For Freshers: Portfolio Over Credentials}

\begin{columns}
    \begin{column}[T]{0.5\linewidth}
\textbf{Engineering freshers:}
\begin{itemize}
\item Deploy at least one app publicly (Streamlit Cloud, HuggingFace Spaces, Vercel)
\item Solve a real problem from your college domain using AI, not a toy dataset
\item Target ``AI-augmented'' roles in your engineering discipline first (AI-enabled civil engineer, AI-enabled QA), not pure ML roles
\item Open-source contributions count heavily at interviews
\end{itemize}
    \end{column}
    \begin{column}[T]{0.5\linewidth}
\textbf{Non-engineering freshers:}
\begin{itemize}
\item Target explicitly: AI Trainer, Prompt Engineer, AI Content Reviewer, AI Policy Analyst
\item Build evidence: case studies showing AI used in your domain work
\item Network via LinkedIn with people already in your target role
\item Look at companies like Scale AI, Outlier, and startups building AI for your sector
\item A short DeepLearning.AI course puts you ahead of 80\% of non-tech applicants
\end{itemize}
    \end{column}
\end{columns}

\vspace{0.1cm}
{\tiny (Ref: DeepLearning.AI Short Courses, free; HuggingFace Spaces, free hosting)}

\end{frame}


%%%%%%%%%%%%%%%%%%%%%%%%%%%%%%%%%%%%%%%%%%%%%%%%%%%%%%%%%%%
\begin{frame}[fragile]\frametitle{For Leads \& Architects: Strategic, Not Reactive}

\begin{columns}
    \begin{column}[T]{0.5\linewidth}
\textbf{Leads:}
\begin{itemize}
\item Dedicate 2 hours/week to structured AI reading (Chip Huyen's blog, The Batch newsletter)
\item Run \textbf{one LLM prototype} on a real team problem in the next 30 days, even a rough one
\item Use AI to amplify your existing output: automate your status reports, code reviews, PRDs
\item Become the person who \textit{evaluates} AI tools critically, not just adopts them blindly
\end{itemize}
    \end{column}
    \begin{column}[T]{0.5\linewidth}
\textbf{Architects:}
\begin{itemize}
\item Adopt evaluation-first design: define how you will measure AI system quality before you build
\item Study RAG vs. fine-tuning vs. prompt patterns, Chip Huyen's ``Designing ML Systems'' is the starting point
\item Design for \textit{model agnosticism}: abstract the LLM layer so you can swap vendors
\item Join LLM engineering communities: Latent Space Discord, MLOps Community
\end{itemize}
    \end{column}
\end{columns}

\vspace{0.1cm}
{\tiny (Ref: Chip Huyen, ``Designing ML Systems'', O'Reilly 2022; Full Stack Deep Learning LLM Bootcamp, free online)}

\end{frame}


%%%%%%%%%%%%%%%%%%%%%%%%%%%%%%%%%%%%%%%%%%%%%%%%%%%%%%%%%%%
\begin{frame}[fragile]\frametitle{For Managers: AI ROI Without AI Theater}

\begin{columns}
    \begin{column}[T]{0.5\linewidth}
\textbf{Short-term actions:}
\begin{itemize}
\item Identify \textbf{2--3 high-ROI use cases} specific to your team's workflows, not generic chatbots
\item Run a 4-week pilot with a real metric (time saved, error rate, output volume)
\item Create a learning budget: structured courses, not just tool subscriptions
\item Designate one team member as the internal AI champion
\end{itemize}
    \end{column}
    \begin{column}[T]{0.5\linewidth}
\textbf{Cultural shift:}
\begin{itemize}
\item Communicate clearly: AI augments the team, it does not replace headcount overnight
\item Reward early experimenters, even when experiments fail
\item Partner with or hire an FDE-type person who can bridge AI and your domain
\item Benchmark against AI-forward peers quarterly; measure outcomes, not adoption numbers
\end{itemize}
    \end{column}
\end{columns}

\vspace{0.2cm}
{\small Key insight: \textbf{the best-managed AI teams are not the ones with the most tools, they are the ones with the clearest use cases and honest measurement.}}

{\tiny (Ref: HBR ``AI Transformation Playbook'' 2024; McKinsey ``The State of AI'' 2024)}
\end{frame}


%%%%%%%%%%%%%%%%%%%%%%%%%%%%%%%%%%%%%%%%%%%%%%%%%%%%%%%%%%%
\begin{frame}[fragile]\frametitle{For Everyone: The AI Literacy Drive}

\begin{columns}
    \begin{column}[T]{0.55\linewidth}
\textbf{Not all of us need to become AI engineers.}

\vspace{0.2cm}
Every professional, in every role, needs \textbf{AI Literacy}, understanding what AI can and cannot do, how to use it responsibly, and how to evaluate its outputs critically.

\vspace{0.2cm}
\textbf{Free, no-code programmes to start with:}
\begin{itemize}
\item \textbf{YUVA AI for ALL}: Future Skills Prime etc (free, Indian languages)
\item \textbf{AI For Everyone}, Andrew Ng, Coursera (free audit)
% \item \textbf{Elements of AI}, University of Helsinki + Reaktor (free)
% \item \textbf{Google AI Essentials}, Coursera (free audit)
% \item \textbf{Microsoft AI Skills Initiative}, learn.microsoft.com
\item \textbf{AI For Bharat / UIWA}, IIT Madras / Govt of India initiatives
\item \textbf{AI4All}, ai4all.org, focused on underrepresented groups
\end{itemize}
    \end{column}
    \begin{column}[T]{0.45\linewidth}
\textbf{Minimum baseline for 2026:}

\begin{itemize}
\item Know what a Large Language Model does and does not do
\item Be able to write a useful prompt for your own work task
\item Understand AI hallucination and how to verify AI output
\item Know when \textit{not} to trust AI
\end{itemize}

\vspace{0.3cm}
{\small Organizations should treat AI literacy the same way they treat data privacy training: \textbf{mandatory, periodic, and role-adapted}.}
    \end{column}
\end{columns}

{\tiny (Ref: ai4all.org; elementsofai.com; learn.microsoft.com/ai; UIWA, Udyami India with AI, Ministry of Skill Development)}

\end{frame}
